% 第 50 章　量子理论的第三条路：路径积分（选读）
\chapter{量子理论的第三条路：路径积分（选读）}

走到这里，你已经见过量子力学的两套行头：薛定谔的波函数语言（第 44 章），一个态随时间连续演化；海森堡的矩阵语言（第 46 章），可观测量变成表格一样的算符。两套语言外观天差地别，算出的预言却逐位相同。这件事本身就是一条线索：也许它们是同一个更深结构的两种投影视角。本章讲第三种写法——费曼的路径积分。它不用波函数，也不用矩阵，只用一个朴素得近乎任性的想法：从起点到终点，把每一条可能的路径都算上。

本章是选读，跳过它不影响后面任何一章。但它有两个别处给不了的好处：其一，它把“量子世界如何长出经典世界”这件事照得最清楚——你将亲眼看到牛顿力学、费马原理、最小作用量原理是怎样从一条更基本的规则里“长”出来的；其二，它是通向量子场论的主路（第 49 章的结尾已经远远望见了它），今天粒子物理的标准计算几乎都用这种语言记账。另外要事先声明一件事：路径积分是一种数学模型和计算规则，不是一组新的实验事实。它做出的每一个可检验预言，都与薛定谔、海森堡的表述完全一致。下面凡是出现“粒子如何如何”的句子，请始终记住这个分层。本章的路线是：先把双缝实验里最刺眼的那个事实按在桌上，再从声音里借来直觉，然后发明费曼的箭头规则，最后用它同时解开双缝之谜和“自然为何选最省事的路”这个老悬念。

\step{反常识开场}

开场还是那台双缝装置，但这一次我们只盯实验事实，不许用任何理论语言粉饰。

电子枪发射电子，穿过挡板上的两条细缝，打在后面的探测屏上，每个电子留下一个亮点。把电子流调到极弱，弱到平均来说装置里同一时刻至多只有一个电子，然后耐心地等。最初几百个点杂乱无章；几千个点之后，隐约露出深浅相间的条纹；几十万个点之后，条纹清晰得像印上去的：一系列等间隔的亮带和暗带，暗带的位置干干净净，一个点都没有。这类“一次只来一个、最后拼出干涉图样”的实验，已经被不同实验室用电子、中子、乃至整个分子反复做过，结果都一样。

现在做关键一步：遮住其中一条缝。直觉说，无非是电子少了一半，图案淡一些而已。事实不是这样。暗带消失了。那些原本几十万个电子里没有一个落下的位置，现在开始不断有电子落下。

把两件事并排摆好，矛盾就无处可躲。对暗带上的某个位置来说：缝一单独开着，电子到得了；缝二单独开着，电子也到得了；两条缝一起开，电子反而永远到不了。而且每个电子都是单独通过装置的，没有同伴可以商量，条纹是几十万个互不相识的电子各自“投票”投出来的。一个孤零零的电子，怎么可能“知道”远处那条它根本没经过的缝是开着还是关着？

先把这个事实按住。它是实验事实，任何理论都必须吞下它。

\step{先猜一猜}

在给出新规则之前，先把你自己的直觉逼出来晒一晒。

猜测一：电子会不会在挡板那里分裂成两半，各走一条缝，到屏前再合起来？如果这个图像是真的，那么在每条缝后面各放一个探测器，应该有机会各抓到“半个电子”。你觉得实验会抓到半个电子，还是每次都抓到一个完整的？

猜测二：有人说，遮住一条缝后暗带变亮，是因为电子“改变了路线”。那么把缝开得更窄，窄到只比电子的德布罗意波长大一点，你觉得“电子经过缝时走的是一条确定的细线”这件事，是更容易说清，还是更难说清？

猜测三：我们以前算两个互斥事件的概率，用的是加法：走缝一到达某处的概率，加上走缝二到达该处的概率。双缝的暗带似乎在嘲笑这条规则。你觉得要修补的是“概率的加法”，还是“相加的对象”？也就是说：是两个概率算错了，还是根本就不该拿概率来相加？

三题的答案本章都会给。先剧透第三题，因为它是全章的钥匙：要改的不是概率，而是相加的对象——量子力学里相加的从来不是概率，而是一种叫作概率幅的东西。

\step{动手实验}

本章的现象在微观尺度，但“两列波动按路程差决定加强还是抵消”这个核心机制，可以在客厅里直接听到。

\begin{experiment}[用两部手机找到“安静点”]
材料：两部手机（或两个便携音箱），一段单频率音频。上网搜索“1 kHz 正弦波”，两部手机播放同一个文件。

第一步，把两部手机并排放在桌上，同时播放，走到房间另一头，音量调到刚好能听清。

第二步，保持两部手机同音量，把它们左右分开约一米，都正对着你。现在捂住一只耳朵，用另一只耳朵听，同时沿与两部手机连线平行的方向慢慢横向移动。

第三步，你会听到音量周期性地变响、变轻，最轻的位置几乎安静。相邻两个“安静点”的间距大约是十七厘米——这正是 1 kHz 声波波长三十四厘米的一半。用尺子量一量，看数量级对不对。

第四步，把两部手机的间距拉开到两米，重走一遍，你会发现“安静点”变得更密。这一步值得记住：两个源的间距越大，干涉条纹越细密。它和双缝实验里“缝距越大条纹越密”是同一条几何规律。

第五步，把其中一部手机静音。那些“安静点”立刻消失，声音变得均匀。
\end{experiment}

解释是第 17、18 章就讲过的：从两部手机到你耳朵的路程一般不相等，两列波到达时振动的步调错开一个相位。路程差恰为波长的整数倍时，两列波同起同落，声音加强；路程差恰为半个波长时，一列往上推一列往下拉，恰好抵消，耳朵附近几乎安静。第五步就是本章开场实验的声音版：两个源同开时听不到声音的位置，关掉一个源反而听得到了。同一个机制还在肥皂泡的彩色条纹、CD 背面的彩虹、降噪耳机里上班（第 35 章），它是经典波动世界最勤劳的员工之一。

这个类比像什么：它像双缝实验里的叠加和相位——路程差决定相位差，相位差决定加强还是抵消，这套数学和光的双缝、电子的双缝是同一个。它不像什么：声波是空气真实密度的起伏，可以直接听、可以用麦克风测；而电子身上“波动”的那个量是概率幅，听不见看不着，只能通过“概率”间接露面。声波相消处能量并没有消失，只是去了别处；量子的暗纹处，是概率幅抵消，这一点后面会看清。

\step{旧模型遇到麻烦}

现在把已有的理论工具都摆上桌，看看它们各自在哪里卡壳。

第一个麻烦：概率加法失灵。初中概率就教过，互斥事件的概率相加。暗纹的存在直接违反这条规则：两条缝各自单独开时，该处的概率都不为零；一起开，概率反而是零。注意，概率论本身没有错。概率能相加的前提是事件互斥而且可以区分。而在双缝实验里，只要你不放探测器，“经缝一到达此处”和“经缝二到达此处”就是同一件观察结果——屏上同一个亮点——的两个无法区分的来历。无法区分来历的事件，不许直接加概率。这是规则层面的事，不是算错了数。

第二个麻烦：薛定谔语言能算，讲起来却别扭。波函数过两条缝、发生干涉、然后在屏上某处“坍缩”出一个点——这套叙述能给出正确答案，但它回答的是“算什么”，没有正面回答另一个问题：为什么这一切到了宏观尺度，就退化成“物体只走一条确定的路径”？波函数弥漫整个空间，台球却沿一条漂亮的抛物线飞行，两者之间那道桥在哪里？

第三个麻烦来自经典物理那边。第 13 章讲过两件像巧合一样的事：光在两点之间走花时间最短的路线（费马原理，第 33、35 章用过它解释折射）；力学系统走的路线使作用量取驻值（最小作用量原理，第 14 章用过它）。两个原理的口吻都像是“自然比较了所有可能的路线，然后挑了一条”。当时这只是一个好用的数学事实：自然凭什么“比较”？它又没有计算机。

这套“比较所有路线”的口吻在日常生活中其实并不陌生。海滩救生员发现有人溺水，从救生亭到溺水者，他不会直线冲过去：沙滩上跑得快，水里游得慢，最省时间的办法是在沙滩上多跑一段再入水，路线在入水点折一个角。光从空气进入水中时的折射，算的是同一类账：光在空气里快、在水里慢，实际光路在界面处弯折的那个角，恰好让总耗时最短。导航软件给你的通勤路线排序时也一样：它给每条候选路线估一个总时间，再把最省的几条排在前面。三个例子里都有一个“沿路线逐段累加的总量”，和最省者胜出的规则。现在把经典和量子两边的麻烦叠在一起，沟的形状就清楚了：一边（量子）是所有来历都参与，一边（经典）是只有一条路线幸存。这条沟上必须有桥，而且桥的两端是同一件事。

顺带记录一个直觉反例，它在后面还要再用：直觉说“通道越多，到达某处的机会越大”。双缝暗纹已经给了它一耳光——两条通道加出一个零。这个反例提醒我们：修补方案必须允许“加一个正贡献，结果反而变小”。普通的正数加法做不到这一点，会转向的量才行。

\step{发明新概念}

1940 年代，费曼把第 13 章那个“自然比较所有路径”的口吻当了真，但做了一次关键的翻转：自然不是比较之后挑一条，而是每一条都计入账本，只是计账的方式会让绝大多数路径互相作废。规则只有三条：

规则一：从起点到终点的每一条可能路径，都分配一个小箭头，所有箭头长度相同。

规则二：每个箭头的指向由这条路径的作用量决定。作用量是沿路径积累的一个量（第 14 章定义过它：动能减势能，再随时间累积）。路径的作用量每积累一个 \(\hbar\)（约化普朗克常数），箭头恰好转过一整圈。作用量大的路径，箭头转得飞快。

规则三：把所有路径的箭头首尾相接地加起来，得到一个合箭头。合箭头长度的平方，正比于“从起点到达终点”这件事发生的概率。

这套规则不是凭空掉下来的。1933 年狄拉克写过一篇短文，提出拉格朗日量——经典力学里那个“动能减势能”——应该在量子理论中扮演某种角色，但没有说清具体怎么扮演。费曼读到了这篇短文，把这个暗示做成了完整的理论：在他的博士论文和 1948 年正式发表的论文里，“每条路径按作用量配一个相位、再全部求和”第一次被写成可以动手计算的工具。一个被搁置了十五年的暗示，长成了量子力学的第三条路。

\begin{center}
\begin{tikzpicture}[scale=1]
  \fill (0,0) circle (1.5pt) node[left=2pt] {源};
  \draw[thick] (2,-1.3) -- (2,-0.32);
  \draw[thick] (2,0.32) -- (2,1.3);
  \draw[thick] (5,-1.3) -- (5,1.3) node[above] {屏};
  \draw (0,0) -- (2,0.5) -- (5,-0.2);
  \draw (0,0) -- (2,-0.5) -- (5,-0.2);
  \fill (5,-0.2) circle (1pt) node[right=2pt] {$x$};
  \node at (1,-0.95) {缝一、缝二各放行一族路径};
\end{tikzpicture}
\end{center}

三条规则说完了，立刻做层次声明，这一章里这件事做多少次都不嫌多。这三条是数学模型的规则，不是实验看到的过程。没有任何实验看到电子“真的走遍所有路径”；“粒子同时走了每一条路”这句话超出了实验能够断言的范围，本书不把它当作事实。反过来，也没有任何实验支持“粒子其实只走了一条确定的细线，只是没告诉我们”。实验能断言的只有一件事：按这三条规则算出的概率，与屏幕上点的分布严丝合缝。至于中间过程，理论部件的私生活，实验无可奉告。

箭头这个比喻也要声明边界。它像音箱实验里的振动：每条路径就像一个声源贡献的一列振动，转多少角度由作用量扮演“路程”的角色，同向加强、反向抵消，数学上完全同构。它不像什么：箭头不是空间中真实的振动，不是电子甩出的尾巴，它只是复数的一种图示（数学桥层会把它写成一个数）；单个箭头也永远测不到，能测的只有合箭头平方给出的概率。

再把导航软件请回来做一次对照。它像本章规则的地方：两者都给每条候选路线算一个数，而且这个数都是沿路线逐段累加出来的——导航累加时间，我们累加作用量。它不像的地方有两处，都很要害：导航算完就把第一名以外的路线全部扔掉，量子规则却把每一条都保留到最后，谁也得不到“被选中”的待遇；导航累加的是只会变大的正数，绕路永远更慢，而量子累加的是会转向的箭头，绕一段路可能让合箭头变长，也可能让它彻底归零。记住这两处差别——“为什么宏观世界只剩一条路径”的答案，就全藏在里面。

\begin{center}
\begin{tikzpicture}[>=Stealth,scale=0.95]
  \draw[->] (0,0) -- (1,0.07);
  \draw[->] (1,0.07) -- (2,-0.02);
  \draw[->] (2,-0.02) -- (3,0.08);
  \draw[->] (3,0.08) -- (4,0.1);
  \draw[->,very thick] (0,0) -- (4,0.1);
  \node at (2,-0.55) {驻值路径附近：箭头几乎同向，合箭头很长};
  \draw[->] (6.4,0) -- (7.3,0.6);
  \draw[->] (7.3,0.6) -- (7.55,-0.35);
  \draw[->] (7.55,-0.35) -- (8.65,0.1);
  \draw[->] (8.65,0.1) -- (8.25,1.0);
  \draw[->] (8.25,1.0) -- (8.5,0.25);
  \draw[->,very thick] (6.4,0) -- (8.5,0.25);
  \node at (7.5,-0.9) {远离驻值路径：箭头打转，合箭头很短};
\end{tikzpicture}
\end{center}

用这三条规则重新看双缝：过缝一的所有路径的箭头加出一个合箭头，记作 \(A_1\)；过缝二的加出 \(A_2\)。屏上某处的总概率幅是 \(A_1+A_2\)。在暗纹位置，\(A_1\) 与 \(A_2\) 恰好等长反向，加起来是零——两条各自畅通的路，合出一片黑暗。这正是概率加法失灵、而箭头加法不失灵的原因。

\step{一句话模型}

\begin{oneline}
从起点到终点，给每一条可能的路径配一个长度相同、方向由作用量决定的箭头（每积累一个 \(\hbar\) 的作用量箭头转一整圈），再把所有箭头加起来；合箭头长度的平方给出事件发生的概率。经典路径并不特殊在“被选中”，而特殊在它附近的箭头几乎指向同一个方向，不互相抵消。
\end{oneline}

注意这句话把两件事一次说清：量子层面“所有路径都计入”，经典层面“只有一条路径显眼”，它们不是两个世界，而是同一条规则在两种账目规模下的两种长相。

\step{数学翻译器}

先把第 14 章的定义搬过来：一条路径的作用量是
\[
S=\int L\,\mathrm{d}t,\qquad L=T-V,
\]
即“动能减势能”沿时间累积。接着把三条箭头规则写成一个公式。每条路径贡献的概率幅正比于
\[
\mathrm{e}^{\mathrm{i}S/\hbar},
\]
这是一个长度为 \(1\)、转角为 \(S/\hbar\) 的复数——正是那个箭头。从起点到终点的总概率幅，是所有路径的贡献之和：
\[
K=\sum_{\text{所有路径}}\mathrm{e}^{\mathrm{i}S/\hbar},
\qquad\text{概率}\propto |K|^2 .
\]

公式出现了，按老规矩先用特殊情形检查它。

检查一：只有两条路径，且 \(S_1=S_2\)。两个箭头同向，合箭头长度是单个的两倍，概率是 \(4\) 倍——不是 \(2\) 倍。概率幅相加后再平方，自然长出干涉项，亮纹比“两缝概率简单相加”更亮。公式没有毛病，这正对应双缝亮纹。

检查二：两条路径，且 \(S_1-S_2=\pi\hbar\)，即箭头恰好转过半圈、首尾相抵。合箭头为零，概率为零。两个都不为零的贡献加出零，这在普通概率里是丑闻，在这个公式里是顺理成章的一行计算。公式再次通过检查：这正对应暗纹。

检查三：令 \(\hbar\) 趋于零，或者等价地，让路径的作用量远大于 \(\hbar\)。此时路径只要偏离一点点，\(S/\hbar\) 就改变成千上万圈，箭头疯狂打转，邻近路径彼此抵消得干干净净。唯一幸免的是驻值路径——作用量取驻值意味着偏离它一阶之内作用量几乎不变，所以它附近的一小簇路径箭头方向几乎一致，集体幸存，加出一根长箭头。第 13 章的“最小作用量原理”就此显形：它从来不是一条独立的自然律法，而是这条求和规则在 \(S\gg\hbar\) 时的近似长相。还要顺手纠正一个历史口气：真正关键的是“驻值”而不是“最小”，第 13 章埋下的这个伏笔在此回收。

还有一个几何图像值得留下。把起点和终点固定，画出一大堆候选路径：直线、弧线、扭来扭去的折线。驻值路径附近的一小簇，箭头几乎同向，合起来是一根长箭头；离它越远，相邻路径的作用量差越大，箭头越转越乱，合箭头缩成一团。真正对概率有发言权的，只是驻值路径周围一条窄窄的“管子”，管子的粗细由 \(\hbar\) 与作用量的比值决定。电子的管子有它的德布罗意波长那么宽，双缝装置里刚好容得下两条这样的管子；台球的管子细到不可思议，于是宏观世界看上去就只剩一条路径。

\begin{center}
\begin{tikzpicture}[scale=1]
  \fill (0,0) circle (1.5pt) node[left=2pt] {起点};
  \fill (6,0) circle (1.5pt) node[right=2pt] {终点};
  \draw[very thick] (0,0) -- (6,0);
  \draw (0,0) .. controls (2,0.5) and (4,0.5) .. (6,0);
  \draw (0,0) .. controls (2,-0.5) and (4,-0.5) .. (6,0);
  \draw (0,0) .. controls (1.5,1.4) and (4.5,1.4) .. (6,0);
  \draw (0,0) .. controls (1.5,-1.4) and (4.5,1.2) .. (6,0);
  \node at (3,0.85) {近处：箭头几乎同向};
  \node at (3,-1.55) {远处：箭头互相打转、彼此作废};
\end{tikzpicture}
\end{center}

\begin{mathbridge}
为什么偏偏是“驻值”路径幸存，可以说得更细。把路径想成可以用一个参数 \(a\) 连续扭动的曲线族，作用量随之变成普通函数 \(S(a)\)。相位 \(S(a)/\hbar\) 随 \(a\) 的变化率是 \(S'(a)/\hbar\)。当 \(\hbar\) 很小，这个变化率极大：参数差一丁点儿，相位就差很多圈，箭头两两反向，求和抵消。唯独在 \(S'(a)=0\) 的点附近，变化率先穿过零，相位在一段不小的参数范围内几乎不动，这一段路径的箭头几乎同向，贡献相干叠加。这就是数学里的“驻相近似”，它是“经典力学是量子力学的近似”这句话最精确的通俗版本。

概率幅还有一条漂亮的组合律：从甲地到丙地、途经某个中间时刻的概率幅，等于“先对所有可能的中间地点乙求和，把甲乙段的概率幅乘上乙丙段的概率幅”。乘法对应箭头接力（转角相加，正是作用量可以逐段累加的反映），求和对应“中间可以在任何地方”。把这条组合律反复使用、把时间切成无穷多小段，就是路径积分的严格定义——它和定积分“切小条再求和”的构造一脉相承，只是被求和的对象从一条曲线下的面积换成了所有曲线本身。
\end{mathbridge}

现在做一个带真实数据的估算，看看“微观所有路径、宏观一条路径”的分界线画在哪里。

先算电子。取一颗动能为 100 eV 的电子（电子显微镜里的典型量级），它的德布罗意波长约 \(1.2\times10^{-10}\) m，速度约 \(6\times10^{6}\) m/s。飞行一米，作用量约为动能乘时间：\(S\sim 1.6\times10^{-17}\times1.7\times10^{-7}\approx3\times10^{-24}\) J·s，即 \(S/\hbar\approx3\times10^{10}\)。已经是天文数字，但还没大到无药可救：两条路径的作用量差取决于路程差，相位差为 \(\pi\) 只需路程差半个波长，约 \(6\times10^{-11}\) m，比原子还小。缝间距做成微米量级，干涉条纹的间距在屏上约 \(10^{-4}\) m，放大后清晰可测。所以电子能“享受”多路径相干的待遇。

再算一粒沙子。质量 \(10^{-6}\) kg、速度 1 m/s、走一米，作用量约 \(5\times10^{-7}\) J·s，于是 \(S/\hbar\approx5\times10^{27}\)。要看到它的干涉，缝间距得和它的德布罗意波长（约 \(7\times10^{-28}\) m）同量级——别说刻缝，这比质子直径还小十几个数量级。对宏观物体，偏离经典路径一万亿亿分之一米，箭头就已转过无数圈，抵消得片甲不留。经典路径不是被自然挑中的，而是唯一没有在求和内战中阵亡的。

\begin{challenge}
这套求和规则凭什么和薛定谔方程给出同样的预言？费曼原始论文的思路是把时间切成 \(N\) 小段，每段极短，路径在每一段上近似直线，于是“所有路径求和”变成 \(N\) 重普通积分的极限；由此可以证明，概率幅随时间的演化恰好满足薛定谔方程。另一个深刻的玩法：把时间换成“虚时间” \(\tau=\mathrm{i}t\)，因子 \(\mathrm{e}^{\mathrm{i}S/\hbar}\) 就变成衰减因子 \(\mathrm{e}^{-S_E/\hbar}\)，剧烈抵消变成温和加权，整个路径积分的模样立刻酷似统计物理里的配分函数（第 31 章）——量子力学与统计力学在公式层面是近亲。这也是为什么路径积分今天成了量子场论的标准语言：在超级计算机上按这个权重对场的大量位形抽样求和（格点量子色动力学），可以从第一性原理算出质子质量，与实验的符合达到百分之一量级。有相互作用时，把相互作用项逐阶展开，求和就表现为一整套带顶点和连线的图——费曼图；第 49 章见到的粒子产生与湮灭的图样，正是它们的速写。
\end{challenge}

\step{模型的边界}

路径积分威力巨大，边界也要划清楚，一共五条。

边界一：它与旧表述严格等价，不是新物理。对非相对论点粒子的任何可检验问题，路径积分、薛定谔方程、海森堡矩阵给出完全相同的预言。三者是同一数学模型的三套账本，选哪套取决于问题：算散射、算场的量子化，路径积分常常最省力；算氢原子能级（第 52 章），薛定谔方程更直接。谁都不是“更真”的那个。

边界二：路径不是旅行日记。求和里的每一条路径只是一个求和指标，和定积分里的小矩形一样，是计算部件。把某条路径指认为“电子真实走的那条路”，就犯了把数学部件当实验事实的错误——这和第 43 章不把电子画成沿确定轨道运行的小球是同一条纪律。

边界三：哪条路信息改变规则本身。一旦装置安排使得路径原则上可以区分——哪怕数据存着不看——相加的对象就从概率幅退化为概率，干涉条纹消失。要强调：这不是仪器太粗糙把电子撞乱了；哪怕扰动任意小、只要路径在原理上可区分，条纹照样消失。原因写在量子的加法规则里，不写在仪器的工艺水平里。

边界四：它没有解决测量问题。路径积分是量子力学的重新表述，不是测量理论的答案。“为什么这一次落在这个点而不是那个点”，它和波函数语言一样只能给出概率，不能给出更多。退相干能解释开放系统里条纹为什么消失，但退相干并不等于回答了全部测量问题——这句话对三套语言同样成立。

边界五：它有框架限制。本章的求和是“固定一个粒子、对它的路径求和”。第 49 章已经看到，高能下粒子数目会变，那时必须升级成对场的一切位形求和。另外，从拉格朗日量的路径语言切回哈密顿量的状态语言需要条件，哈密顿量并不无条件等于总能量——这层小心在第 15 章提过，在这里同样有效。

最后还要提防一种流行的误用：“量子力学说粒子走所有路径，所以一切皆有可能。”账本里确实记着稀奇古怪的路径——绕月亮一圈再落到屏上的路径也在其中——但它们的箭头几乎全部彼此作废，对概率的贡献小到无法形容。“写进账本”和“产生可见后果”之间，隔着干涉相消这道天堑。路径积分扩大的是计算的视野，不是奇迹的配额。

\step{回头解释开场现象}

现在回到开场的悖论，逐条结账。

暗纹处为什么一个点也没有？过缝一的所有路径加出合箭头 \(A_1\)，过缝二的加出 \(A_2\)，在该处 \(A_1\) 与 \(A_2\) 恰好等长反向，总概率幅为零，概率为零。

遮住一条缝，为什么那里反而有电子落下？遮缝删掉了整组 \(A_2\)，\(A_1\) 单独存在，概率不再为零。电子没有分裂，没有商量，也没有以任何方式“感知”另一条缝；改变的只是我们计算概率幅时的求和范围。装置变了，事件集合变了，概率分布当然跟着变——怪的不是电子，是我们不肯放弃的概率加法。

条纹的样子也能顺手解释。从暗纹移到相邻亮纹，需要的是两组路径的作用量差改变 \(\pi\hbar\)，也就是路程差改变半个德布罗意波长。于是缝距越大，同样路程差对应的屏幕位置越近，条纹就越细密——这正是音箱实验第四步里“源距越大、安静点越密”的量子版本。声音、光、电子，三层世界里刻着同一条几何规律，只是每层世界里“转向的那个量”换成了不同的主角：声压、电场、概率幅。

开场三道猜测的答案也一并交卷。猜测一：每条缝后放探测器，抓到的永远是完整电子，半个电子从来不存在；但代价是路径变成可区分，条纹消失。猜测二：缝越窄，过缝位置定得越准，按不确定性关系，横向动量就越不准，条纹反而摊得越宽——想说清“一条确定的细线”反而更难，这不是仪器限制，是量子规则本身。猜测三：要改的是相加的对象——概率不可直接相加，概率幅可以；概率幅是可正可负可转向的量，“加一个贡献结果变小”不再是悖论，而是它的家常便饭。

还有一层安静但重要的回报：第 13 章“自然选择最省事的路”那个拟人化口吻，终于可以退休了。自然没有比较，也没有选择；它只是让每一条路径都投票，而在宏观尺度上，只有驻值路径附近的选票方向一致，其余全部作废。最小作用量原理是统计结果，不是立法条文。

\step{脑洞实验与挑战题}

\begin{thought}[把缝开到无穷多]
双缝装置开三条缝，概率幅就是三族路径的箭头相加；开十条、一千条、密密麻麻直到挡板千疮百孔……最后干脆把挡板撤掉。按同一逻辑推到底：自由粒子的传播本身就是对自由空间中所有路径的求和。这不是数学杂技，正是路径积分的诞生思路——双缝不是特例，双缝是把一般规则裁剪到只剩两族路径的极简版。挡板的全部作用，只是“删除若干求和项”。

再推一个极端尺度。想象一个 \(\hbar\) 不是 \(10^{-34}\) J·s 而是 1 J·s 的世界。一颗台球的作用量大约就是几个 \(\hbar\)，偏离直线的路径不再互相抵消干净：台球从球桌这头到那头会散开成一朵概率云，途中立一块开双缝的挡板，它落点的分布会出现干涉条纹。这个思想实验的教益是：量子与经典的边界不在“尺寸大小”，而在作用量与 \(\hbar\) 的比值。大物体之所以经典，是因为它们的账本上 \(\hbar\) 太小，小到任何路径分歧都立刻被内战抹平。
\end{thought}

这条路通向哪里？把“对路径求和”升级为“对场求和”，就得到量子场论的标准表述，第 49 章结尾的图景正是用它算账的；再往前，把求和对象换成时空几何本身，就是量子引力研究正在试探的方向（第 61 章的地图上标着它）。一条从双缝出发的规则，一路通到物理学最前沿的工地。

收尾时，用全书一直在追问的三件事给本章对一次账。系统此刻是什么状态？路径积分不换答案，只换视角：不必追问粒子“此刻在哪里”，要问从初态出发、到达每个可能结果的概率幅是多少，而这个幅是所有路径箭头的和。什么规律决定状态如何变化？作用量。它给每条路径定相位，是这条路上唯一的立法者，经典力学只是它的近似剪影。我们怎样通过测量知道它？仍然只有概率：合箭头长度的平方，一次实验一个亮点，条纹永远靠积累。三套语言，同一份答卷——这本身就是“它们背后是同一个实在”最有力的证词。

\begin{puzzles}
\item 把双缝改成三缝。双缝时的某个暗纹位置，三条缝全开时一定还是暗的吗？（思路提示：暗纹要求所有合箭头加起来为零。两个等长箭头要抵消只能反向；三个箭头要抵消，只需能首尾相接围成一个三角形——条件宽松得多。）
\item 在缝一后面贴一片极薄的玻璃，电子穿过玻璃时作用量略有改变，相当于 \(A_1\) 整体多转了一个角度，\(A_2\) 不动。干涉条纹会怎样？（思路提示：暗纹位置取决于两个箭头的相对角度；相对角度变了，满足“等长反向”的屏幕位置就搬家——条纹整体平移，但深浅间距不变。实验中这确实是移动条纹的标准手段。）
\item 路径积分说“对所有路径求和”，为什么双缝问题里只算“过缝一”和“过缝二”两族就够了？撞上挡板的那些路径呢？（思路提示：求和的对象是“从源到达屏上某点”这一事件的全部候选路径；被挡板吸收的路径到不了屏，不构成该事件的候选。改了挡板，就改了事件的定义，求和范围随之改变——这正是遮缝改变概率的算术根源。）
\item 一颗质量 \(10^{-6}\) kg 的沙粒以 1 m/s 运动，想用双缝看到它的干涉条纹，缝间距大约要什么量级？（思路提示：条纹由德布罗意波长决定，\(\lambda=h/(mv)\)。算出波长的数量级，再想想人类能不能刻出这样的缝——这就是宏观世界“只有一条路径”的定量原因。）
\end{puzzles}
